\documentclass[11pt,a4paper]{article}

% ---------- Codificacion / idioma / fuentes (XeLaTeX + tectonic) ----------
\usepackage{fontspec}
\usepackage[spanish,es-noindentfirst]{babel}

% ---------- Geometria y estilo general ----------
\usepackage{geometry}
\geometry{margin=2.4cm}
\usepackage{xcolor}
\usepackage{enumitem}
\usepackage{booktabs}
\usepackage{titlesec}
\usepackage{fancyhdr}
\usepackage{listings}
\usepackage[most]{tcolorbox}
\usepackage{hyperref}
\usepackage{parskip}

% ---------- Colores ----------
\definecolor{brand}{HTML}{1F4E79}
\definecolor{brandlight}{HTML}{EAF1F8}
\definecolor{codebg}{HTML}{F6F8FA}
\definecolor{codeframe}{HTML}{D0D7DE}
\definecolor{kw}{HTML}{CF222E}
\definecolor{str}{HTML}{0A3069}
\definecolor{com}{HTML}{6E7781}
\definecolor{ok}{HTML}{1A7F37}
\definecolor{bad}{HTML}{B42318}

% ---------- Hyperref ----------
\hypersetup{colorlinks=true, linkcolor=brand, urlcolor=brand, citecolor=brand}

% ---------- Titulos ----------
\titleformat{\section}{\Large\bfseries\color{brand}}{\thesection.}{0.6em}{}
\titleformat{\subsection}{\large\bfseries\color{brand!85!black}}{\thesubsection}{0.6em}{}
\titlespacing*{\section}{0pt}{1.4em}{0.6em}

% ---------- Encabezado / pie ----------
\pagestyle{fancy}
\fancyhf{}
\lhead{\small\color{brand}\bfseries Reparación Automática de Software}
\rhead{\small\color{brand} Práctica 1}
\cfoot{\small\thepage}
\renewcommand{\headrulewidth}{0.4pt}

% ---------- Estilo de codigo ----------
\lstdefinestyle{py}{
  language=Python,
  backgroundcolor=\color{codebg},
  basicstyle=\ttfamily\small,
  keywordstyle=\color{kw}\bfseries,
  stringstyle=\color{str},
  commentstyle=\color{com}\itshape,
  numberstyle=\tiny\color{com},
  numbers=left, numbersep=8pt,
  frame=single, rulecolor=\color{codeframe},
  framesep=6pt, xleftmargin=14pt,
  showstringspaces=false,
  breaklines=true, breakatwhitespace=true,
  tabsize=4,
  literate={á}{{\'a}}1 {é}{{\'e}}1 {í}{{\'i}}1 {ó}{{\'o}}1 {ú}{{\'u}}1 {ñ}{{\~n}}1
}
\lstdefinestyle{sh}{
  language=bash,
  backgroundcolor=\color{codebg},
  basicstyle=\ttfamily\small,
  keywordstyle=\color{brand}\bfseries,
  commentstyle=\color{com}\itshape,
  frame=single, rulecolor=\color{codeframe},
  framesep=6pt, xleftmargin=6pt,
  showstringspaces=false, breaklines=true, tabsize=2,
}

% ---------- Cajas ----------
\newtcolorbox{answer}{colback=brandlight, colframe=brand, boxrule=0.8pt,
  arc=3pt, left=8pt, right=8pt, top=6pt, bottom=6pt}

\newtcolorbox{def-err-fail}{colback=white, colframe=bad!70!black, boxrule=0.8pt,
  arc=3pt, left=8pt, right=8pt, top=6pt, bottom=6pt}

% Etiquetas
\newcommand{\lblP}{\textcolor{ok}{\textbf{[PASA]}}}
\newcommand{\lblF}{\textcolor{bad}{\textbf{[FALLA]}}}
\newcommand{\defecto}{\textbf{\textcolor{bad}{Defecto (fault):}}\ }
\newcommand{\error}{\textbf{\textcolor{bad}{Error (infección):}}\ }
\newcommand{\falla}{\textbf{\textcolor{bad}{Falla (failure):}}\ }
\newcommand{\code}[1]{\colorbox{codebg}{\ttfamily\small #1}}

% =====================================================================
\begin{document}

% ---------- Portada ----------
\begin{center}
  {\LARGE\bfseries\color{brand} Práctica 1 --- Debugging y Prueba de Código}\\[0.4em]
  {\large Reparación Automática de Software --- 2026}\\[0.2em]
  {\normalsize Universidad Nacional de Río Cuarto (UNRC)}\\[1.2em]
\end{center}

\begin{answer}
\small\textbf{Resumen.} Este informe resuelve los cuatro ejercicios de la Práctica 1.
Los Ejercicios 1 y 2 documentan la instalación de \emph{The Debugging Book} y los conceptos
centrales de sus capítulos \emph{01\_Intro} y \emph{02\_Observing}. Los Ejercicios 3 y 4
analizan los módulos \code{max.py} y \code{person.py}: para cada test que falla se identifica el
\textbf{defecto}, el \textbf{error} y la \textbf{falla}, y se propone la corrección.
\textbf{Todos los resultados de los tests fueron verificados ejecutando el código} (Python 3.9).
\end{answer}

% ---------- Vocabulario base ----------
\section*{\normalsize\color{brand}Vocabulario base (Cap.\ 1)}
\vspace{-0.4em}
Los tres términos que se usan en todo el informe (del capítulo \emph{Introduction to Debugging}):
\begin{itemize}[leftmargin=1.4em, itemsep=2pt]
  \item \defecto la porción de código escrita incorrectamente (la causa raíz en el texto del programa).
  \item \error un estado interno del programa durante la ejecución que es incorrecto por culpa del defecto (el texto lo llama \emph{infection}). Un defecto puede no producir error si nunca se ejecuta o el estado malo no llega a propagarse.
  \item \falla el comportamiento observable e incorrecto (salida errónea, excepción, resultado equivocado) que percibe quien usa el programa.
\end{itemize}
La cadena causal es: \textbf{Defecto} $\rightarrow$ \textbf{Error} $\rightarrow$ \textbf{Falla}.

\newpage
% =====================================================================
\section{Instalación de \emph{The Debugging Book}}

Pasos ejecutados sobre Ubuntu 22.04+ (x86-64), siguiendo la guía:

\begin{enumerate}[leftmargin=1.6em, itemsep=3pt]
  \item Descargar desde \href{https://www.debuggingbook.org/}{debuggingbook.org} el paquete de
        notebooks: \emph{resources} $\rightarrow$ \emph{All Notebooks (.zip)}.
  \item Descomprimir y aplicar los ajustes previos (Graphviz + fix de \code{requirements.txt}).
\end{enumerate}

\begin{lstlisting}[style=sh]
# 1. Descomprimir el material
unzip debuggingbook-notebooks.zip
cd debuggingbook-notebooks     # carpeta resultante

# 2. Fix requerido por la guia (punto d): en requirements.txt,
#    la linea 30 debe empezar con "showast2" en lugar de "showast".
sed -i '30s/^showast/showast2/' requirements.txt

# 3. Asegurar Graphviz (necesario para los grafos de las notebooks)
dot --version || sudo apt install graphviz

# 4. Crear y activar un entorno virtual aislado
python3 -m venv .venv
source .venv/bin/activate

# 5. Instalar dependencias  (OJO: la guia dice "requirements"
#    pero el comando correcto lleva la extension .txt)
pip install -r requirements.txt

# 6. Levantar las notebooks
jupyter notebook
# Si no abre el navegador: http://localhost:8888/tree
\end{lstlisting}

\begin{answer}\small
\textbf{Notas / correcciones respecto del enunciado:}
\begin{itemize}[leftmargin=1.3em, itemsep=1pt, topsep=2pt]
  \item El fix del punto (d) evita el conflicto de nombres del paquete \code{showast} por \code{showast2}.
  \item El punto (h) dice \code{pip install -r requirements}; el comando correcto es
        \code{pip install -r requirements.txt}.
  \item Conviene crear el \code{venv} \emph{dentro} de la carpeta descomprimida para mantener todo aislado.
\end{itemize}
\end{answer}

% =====================================================================
\section{Capítulos \emph{01\_Intro} y \emph{02\_Observing}}

\subsection*{01 --- Introduction to Debugging}
Ideas centrales completadas:
\begin{itemize}[leftmargin=1.4em, itemsep=2pt]
  \item Distinción \textbf{defecto / error / falla} (ver vocabulario base). Es el marco que se usa en los Ejercicios 3 y 4.
  \item \textbf{El método científico del debugging}: observar la falla, formular una \emph{hipótesis} sobre la causa, predecir una consecuencia observable, experimentar para confirmarla o refutarla, y repetir hasta aislar el defecto. Es un proceso iterativo, no adivinanza.
  \item \textbf{Debugging = búsqueda de la causa raíz}, no sólo enmascarar el síntoma: corregir el defecto, no la falla.
  \item Ejemplo canónico del libro: \code{remove\_html\_markup}, donde un caso con comillas rompe la función y sirve para practicar la formulación de hipótesis.
\end{itemize}

\subsection*{02 --- How to Observe}
Técnicas para observar el estado del programa y localizar el error:
\begin{itemize}[leftmargin=1.4em, itemsep=2pt]
  \item \textbf{Observación con \code{print}/logging}: rápida pero invasiva; útil para un primer sondeo.
  \item \textbf{Depurador interactivo} (\code{pdb}, breakpoints): inspeccionar variables, avanzar paso a paso, examinar el \emph{stack}.
  \item \textbf{Aserciones (\code{assert})}: codifican invariantes; detectan el error \emph{cerca} de su origen, antes de que se propague hasta la falla.
  \item El objetivo es acortar la distancia entre \emph{defecto} y \emph{error observado}: cuanto antes se detecta el estado infectado, más fácil es ubicar la línea culpable.
\end{itemize}

% =====================================================================
\section{Módulo \code{max.py}}

El módulo tiene dos funciones. \code{max\_int} es correcta; el defecto está en \code{max\_int\_seq}.

\subsection{Tests}
Archivo entregable: \code{test\_max.py} (ejecutable con \code{pytest -v}).

\begin{lstlisting}[style=py]
from max import max_int, max_int_seq

# --- max_int: correcta, todos pasan ---
def test_max_int_a_mayor():   assert max_int(5, 3) == 5   # PASA
def test_max_int_b_mayor():   assert max_int(3, 5) == 5   # PASA
def test_max_int_iguales():   assert max_int(4, 4) == 4   # PASA

# --- max_int_seq ---
# Spec asumida (segun el codigo): distinta longitud -> gana la mas larga;
# igual longitud -> gana la que tiene mas maximos elemento-a-elemento.
def test_seq_distinta_longitud():
    assert max_int_seq([1], [2, 3]) == [2, 3]            # PASA
def test_seq_igual_longitud_gana_b():
    assert max_int_seq([1, 2], [3, 4]) == [3, 4]         # PASA (por coincidencia)
def test_seq_igual_longitud_gana_a():
    assert max_int_seq([9, 9], [1, 1]) == [9, 9]         # FALLA
def test_seq_igual_longitud_gana_a_2():
    assert max_int_seq([3, 5], [1, 2]) == [3, 5]         # FALLA
\end{lstlisting}

\textbf{Resultado verificado:} los tests de \code{max\_int} y el de distinta longitud \lblP{}.
El caso en que \code{b} gana legítimamente \lblP{} \emph{por coincidencia} (ver abajo).
Los dos últimos \lblF{}: para \code{[9,9]} vs \code{[1,1]} la función devuelve \code{[1,1]} cuando
debería devolver \code{[9,9]}.

\subsection{Defecto, error y falla}
\begin{def-err-fail}
\defecto en la línea 28, la comparación es
\code{if max\_int(a[i], b[i]) == a:} — compara el máximo (un \emph{int}) contra la \emph{lista} completa \code{a}.
Debería comparar contra el elemento \code{a[i]}.
\smallskip

\error \code{max\_int(a[i], b[i])} devuelve un \code{int}, que nunca es igual a la lista \code{a};
por lo tanto la condición es \emph{siempre falsa}. El contador \code{max\_a} queda en \code{0} y sólo
se incrementa \code{max\_b}, aunque \code{a} tenga todos los máximos. El estado interno de los
contadores es incorrecto.
\smallskip

\falla en el caso de igual longitud la función \emph{siempre} devuelve \code{b}. Con
\code{max\_int\_seq([9,9],[1,1])} retorna \code{[1,1]} en lugar de \code{[9,9]}.
\smallskip

\emph{Coincidencia:} el test con \code{[1,2]} vs \code{[3,4]} pasa porque ahí \code{b}
efectivamente debía ganar; el bug queda oculto (\emph{coincidental correctness}).
\end{def-err-fail}

\subsection{Corrección}
\begin{lstlisting}[style=py]
        for i in range(0, len(a)):
            if max_int(a[i], b[i]) == a[i]:   # FIX: a[i] en lugar de a
                max_a += 1
            else:
                max_b += 1
\end{lstlisting}
Con este cambio, los cuatro tests de \code{max\_int\_seq} pasan (verificado).
Observación adicional: en caso de empate (\code{max\_a == max\_b}) la función devuelve \code{b};
es una decisión de diseño razonable pero conviene documentarla.

% =====================================================================
\section{Módulo \code{person.py}}

Se responden las preguntas embebidas como \code{\#TODO} en el módulo. Archivo de tests entregable:
\code{test\_person.py}. Las fechas se calculan relativas a \emph{hoy} para que los tests sean
deterministas.

% ---------------- __init__ ----------------
\subsection{\code{\_\_init\_\_} --- ¿Es correcta la implementación?}
\begin{answer}
\textbf{Respuesta:} funciona como simple contenedor de datos, pero \textbf{no valida} sus argumentos,
por lo que \emph{permite construir estados inválidos} (p.\ ej.\ una fecha de nacimiento en el futuro,
o \code{studies} que no sea un \code{Studies}). Bajo la especificación ``un \code{Person} no debe
poder estar en un estado inválido'', la implementación es \emph{incompleta}.
\end{answer}

\begin{lstlisting}[style=py]
def test_init_guarda_atributos():           # PASA
    p = Person(date(2000,1,1), "David", Studies.BACHELOR)
    assert p.name == "David" and p.studies == Studies.BACHELOR

def test_init_no_valida_fecha_futura():     # FALLA (segun spec)
    futuro = date(date.today().year + 5, 1, 1)
    p = Person(futuro, "David")
    assert p.birthdate <= date.today()      # el constructor deberia rechazarla
\end{lstlisting}

\begin{def-err-fail}
\defecto ausencia de validación en el constructor.\quad
\error el objeto queda con \code{\_birthdate} en el futuro (estado inválido).\quad
\falla luego \code{age} devuelve una edad negativa y \code{\_\_str\_\_} muestra datos absurdos.
\end{def-err-fail}

\begin{answer}\small\textbf{Corrección sugerida:} validar en el constructor.
\begin{lstlisting}[style=py]
def __init__(self, birth, name, studies=Studies.NONE):
    if not isinstance(birth, date) or birth > date.today():
        raise ValueError("birth debe ser una fecha no futura")
    if not name:
        raise ValueError("name no puede ser vacio")
    if not isinstance(studies, Studies):
        raise TypeError("studies debe ser un Studies")
    self._birthdate = birth; self._name = name; self._studies = studies
\end{lstlisting}
\end{answer}

% ---------------- mutabilidad ----------------
\subsection{¿Sería correcto permitir modificar la fecha de nacimiento y el nombre?}
\begin{answer}
\textbf{Fecha de nacimiento: no.} Es un dato inmutable en la realidad (nunca cambia); permitir
modificarla sólo habilita estados inconsistentes. Por eso el módulo expone \code{birthdate} sólo
como \emph{property} de lectura, sin \emph{setter}: es la decisión correcta.
\smallskip

\textbf{Nombre: es discutible.} Un nombre \emph{puede} cambiar legalmente, así que un \emph{setter}
sería defendible. Pero si el nombre forma parte de la identidad del objeto (se usa en \code{\_\_eq\_\_}
y en un eventual \code{\_\_hash\_\_}), hacerlo mutable rompe invariantes al usar el objeto en
\code{set}/\code{dict}. Mantenerlo inmutable (como está) es la opción más segura.
\end{answer}

% ---------------- age ----------------
\subsection{\code{age} --- ¿Es correcta la implementación?}
\begin{answer}
\textbf{Respuesta: no.} Calcula la edad restando sólo los años y \emph{ignora si el cumpleaños ya
ocurrió} este año. Para quien nació más tarde en el año que la fecha actual, sobreestima la edad en~1.
\end{answer}

\begin{lstlisting}[style=py]
def test_age_cumpleanios_ya_paso():         # PASA
    p = Person(date(date.today().year - 30, 1, 1), "David")  # nacio 1-ene
    assert p.age == 30

def test_age_cumpleanios_no_llego():        # FALLA (off-by-one)
    p = Person(date(date.today().year - 30, 12, 31), "David")  # nacio 31-dic
    assert p.age == 29     # hasta diciembre todavia tiene 29
\end{lstlisting}

\begin{def-err-fail}
\defecto \code{years = today.year - birth.year}, sin ajustar por mes/día.\quad
\error para un nacido el 31-dic, \code{years} vale \code{30} cuando el valor correcto es \code{29}.\quad
\falla \code{age} retorna un número mayor al real (verificado: devuelve \code{30} en vez de \code{29}).
\end{def-err-fail}

\begin{answer}\small\textbf{Corrección:} restar 1 si el cumpleaños de este año aún no llegó.
\begin{lstlisting}[style=py]
today = date.today()
years = today.year - self._birthdate.year - (
    (today.month, today.day) < (self._birthdate.month, self._birthdate.day)
)
return years
\end{lstlisting}
\end{answer}

% ---------------- studies setter ----------------
\subsection{\code{studies.setter} --- ¿Es correcta la implementación?}
\begin{answer}
\textbf{Respuesta: no.} El \emph{setter} asigna directamente sin validar. Esto (a)~acepta valores que
no son \code{Studies}, rompiendo la invariante de tipo, y (b)~permite \emph{bajar} el nivel de estudios,
lo cual no tiene sentido (no se ``des-estudia'' un título).
\end{answer}

\begin{lstlisting}[style=py]
def test_studies_setter_actualiza():        # PASA
    p = Person(date(2000,1,1), "David", Studies.NONE)
    p.studies = Studies.MASTERS
    assert p.studies == Studies.MASTERS

def test_studies_setter_valida_tipo():      # FALLA (sin validacion)
    p = Person(date(2000,1,1), "David", Studies.MASTERS)
    p.studies = "PhD"                        # valor invalido
    assert isinstance(p.studies, Studies)

def test_studies_setter_no_regresa():       # FALLA (permite regresion)
    p = Person(date(2000,1,1), "David", Studies.MASTERS)
    p.studies = Studies.PRIMARY
    assert p.studies == Studies.MASTERS
\end{lstlisting}

\begin{def-err-fail}
\defecto asignación sin validación (\code{self.\_studies = new\_studies}).\quad
\error \code{\_studies} queda con un \code{str} arbitrario o con un nivel menor al previo.\quad
\falla un \code{Person} con nivel \code{MASTERS} pasa a \code{PRIMARY}, o \code{\_studies} deja de ser
un \code{Studies} (verificado: ambas asignaciones se aceptan sin error).
\end{def-err-fail}

\begin{answer}\small\textbf{Corrección:} validar tipo e impedir regresión.
\begin{lstlisting}[style=py]
@studies.setter
def studies(self, new_studies):
    if not isinstance(new_studies, Studies):
        raise TypeError("studies debe ser un Studies")
    if new_studies.value[1] < self._studies.value[1]:
        raise ValueError("el nivel de estudios no puede disminuir")
    self._studies = new_studies
\end{lstlisting}
\end{answer}

% ---------------- __eq__ ----------------
\subsection{\code{\_\_eq\_\_} --- ¿Es correcta la implementación?}
\begin{answer}
\textbf{Respuesta: no}, tiene dos defectos. (1)~Compara los nombres con \code{is not} (identidad de
objetos) en lugar de \code{!=} (igualdad de valor). (2)~Compara \code{age} en vez de \code{\_birthdate},
lo que hace que la igualdad dependa del día en que se ejecute y que dos personas de distinta fecha
pero misma edad se consideren iguales. Además, al definir \code{\_\_eq\_\_} sin \code{\_\_hash\_\_},
la clase queda \emph{no hasheable}.
\end{answer}

\begin{lstlisting}[style=py]
def test_eq_mismos_datos_nombre_interno():  # PASA (por interning)
    a = Person(date(1990,5,5), "David", Studies.NONE)
    b = Person(date(1990,5,5), "David", Studies.NONE)
    assert a == b

def test_eq_nombre_no_interno():            # FALLA (por `is not`)
    n1 = "".join(["D","a","v","i","d"]); n2 = "".join(["D","a","v","i","d"])
    a = Person(date(1990,5,5), n1, Studies.NONE)
    b = Person(date(1990,5,5), n2, Studies.NONE)
    assert a == b

def test_person_es_hasheable():             # FALLA (falta __hash__)
    a = Person(date(1990,5,5), "David", Studies.NONE)
    assert hash(a) is not None
\end{lstlisting}

\begin{def-err-fail}
\defecto \code{if self.\_name is not other.\_name} usa identidad; y se compara \code{age} en vez de \code{\_birthdate}.\quad
\error para dos nombres iguales pero construidos en runtime (no \emph{interned}), \code{is not} da
\code{True} y la función retorna \code{False} prematuramente.\quad
\falla \code{a == b} devuelve \code{False} siendo lógicamente iguales (verificado). Y \code{hash(a)}
lanza \code{TypeError: unhashable type}.
\end{def-err-fail}

\begin{answer}\small\textbf{Corrección:} comparar por valor, usar \code{\_birthdate} y agregar \code{\_\_hash\_\_}.
\begin{lstlisting}[style=py]
def __eq__(self, other):
    if not isinstance(other, Person):
        return False
    if self._name != other._name:            # FIX: != en lugar de `is not`
        return False
    if self._birthdate != other._birthdate:  # FIX: comparar fecha, no age
        return False
    if self._studies != other._studies:
        return False
    return True

def __hash__(self):
    return hash((self._name, self._birthdate, self._studies))
\end{lstlisting}
\end{answer}

% ---------------- max ----------------
\subsection{\code{max(a, b)} --- Implementar y proveer tests}
\begin{answer}
La consigna no fija qué significa ``persona máxima'', así que se define un
\textbf{orden explícito y documentado}: gana quien tiene \emph{mayor nivel de estudios}; a igual nivel,
gana quien tiene \emph{mayor edad} (fecha de nacimiento más temprana); si empatan en todo, se devuelve \code{a}.
\end{answer}

\begin{lstlisting}[style=py]
def max(a: Person, b: Person) -> Person:
    """Persona maxima: mayor nivel de estudios; a igualdad, mayor edad."""
    na, nb = a.studies.value[1], b.studies.value[1]   # nivel ISCED interno
    if na != nb:
        return a if na > nb else b
    if a.birthdate != b.birthdate:
        return a if a.birthdate < b.birthdate else b   # mas edad = fecha menor
    return a
\end{lstlisting}

\begin{lstlisting}[style=py]
def test_max_por_estudios():        # PASA con la implementacion
    a = Person(date(1990,1,1), "David", Studies.DOCTORATE)
    b = Person(date(1990,1,1), "Isaac", Studies.BACHELOR)
    assert max(a, b) is a

def test_max_desempata_por_edad():  # PASA con la implementacion
    a = Person(date(1970,1,1), "David", Studies.MASTERS)   # mas grande
    b = Person(date(2000,1,1), "Isaac", Studies.MASTERS)
    assert max(a, b) is a
\end{lstlisting}
\begin{answer}\small
\textbf{Criterio verificado} ejecutando ambos tests. Si la cátedra prefiere otro orden (p.\ ej.\ sólo
por edad), basta cambiar la clave de comparación; la estructura del código y de los tests se mantiene.
\end{answer}

% =====================================================================
\section*{\normalsize\color{brand}Resumen de defectos encontrados}
\vspace{-0.3em}
\begin{center}
\small
\begin{tabular}{@{}p{2.6cm}p{4.5cm}p{7.2cm}@{}}
\toprule
\textbf{Módulo / función} & \textbf{Defecto} & \textbf{Falla observable} \\
\midrule
\code{max\_int\_seq} & \code{== a} en vez de \code{== a[i]} & Con igual longitud siempre devuelve \code{b}. \\
\code{Person.\_\_init\_\_} & Sin validación de argumentos & Permite fecha futura / tipos inválidos. \\
\code{Person.age} & No ajusta por mes/día & Edad sobreestimada en 1 antes del cumpleaños. \\
\code{Person.studies} (setter) & Sin validación & Acepta tipos inválidos y permite bajar de nivel. \\
\code{Person.\_\_eq\_\_} & \code{is not} en vez de \code{!=}; usa \code{age} & Personas iguales dan \code{!=}; clase no hasheable. \\
\code{Person.max} & No implementada & Lanzaba \code{NotImplementedError}; ya implementada. \\
\bottomrule
\end{tabular}
\end{center}

\end{document}
